\begin{RTL}
در سنتز آنتی بیوتیک سفاماندول، مکانیزم تبدیل مرحله ی آخر که در زیر نمایش داده شده است و با خروج دی اکسید کربن است را بنویسید.
\end{RTL}

\begin{center}
\fbox{%
\begin{minipage}{0.95\linewidth}
\schemestart
\chemname{\CephemNHtwo}{}
\arrow{->[\chemfig{Ph-C(=O)-O-C(=O)-O}]} % anhydride (simplified)
\chemname{\CephemMand}{\textcolor{purple}{Cefamandole}}
\schemestop
\end{minipage}}
\end{center}

\begin{RTL}
\textbf{مکانیزم مرحله آخر (خروج CO2):}
\end{RTL}

\begin{center}
\begin{tabular}{cccc}
\begin{minipage}{0.23\linewidth}\centering
\CephemMand
\par\scriptsize 1) حمله گروه OH به کربونیل بتالاکتام
\end{minipage} &
\begin{minipage}{0.23\linewidth}\centering
\CephemMandTetra
\par\scriptsize 2) تشکیل واسطه تتراهدرال
\end{minipage} &
\begin{minipage}{0.23\linewidth}\centering
\CephemMand
\par\scriptsize 3) باز شدن حلقه بتالاکتام
\end{minipage} &
\begin{minipage}{0.23\linewidth}\centering
\CephemMand
\par\scriptsize 4) انتقال پروتون
\end{minipage}
\end{tabular}

\vspace{0.6em}

\begin{tabular}{ccc}
\begin{minipage}{0.31\linewidth}\centering
\CephemMand
\par\scriptsize 5) ایجاد گروه ترک کننده مناسب
\end{minipage} &
\begin{minipage}{0.31\linewidth}\centering
\CephemMand \;+\; \chemfig{O=C=O}
\par\scriptsize 6) خروج CO2
\end{minipage} &
\begin{minipage}{0.31\linewidth}\centering
\CephemMand
\par\scriptsize 7) تشکیل محصول نهایی (سفاماندول)
\end{minipage}
\end{tabular}
\end{center}